\documentclass[12pt]{article}

\usepackage{fontawesome}
\usepackage{hyperref}
\usepackage{xurl}
\usepackage[tagged, highstructure]{accessibility}
\usepackage{bookmark}
\usepackage{enumitem}

\hypersetup{
    colorlinks=false,
    pdfborder={0 0 0},
    pdftitle={Scripts with Looping},
    pdfauthor={Adrianna Holden-Gouveia},
    pdfsubject={Introduction to Linux - Lab Assignment},
    pdflang={en-US},
    pdfkeywords={lab assignment, Scripts with Looping}
}

% Configure PDF bookmarks for navigation
\bookmarksetup{
    numbered,
    open,
}

% Configure list spacing for better accessibility
\setlist{nosep}


\title{Scripts with Looping}
\author{
        Adrianna Holden-Gouveia \\
        Website: \href{https://aholdengouveia.name}{aholdengouveia.name}\\ 
        \date{\vspace{-5ex}}
        %Email: \href{mailto:admin@aholdengouveia.name}{admin@aholdengouveia.name} \\
%        \faLinkedin{: aholdengouveia} \\
        \faGithub {: aholdengouveia} \\
%        \faTwitter {: aholdengouveia} \\
        }

%S\date{\today}


\begin{document}    

\maketitle

%\begin{abstract}
%This is a template for Linux Administration Lab work
%\end{abstract}
%\tableofcontents
\section*{Accessibility Notice}
This document is also available in HTML format at:

Link: \href{/home/aholdengouveia/aholdengouveia.github.io/IntroLinux/labs/shellloop.html}{\textbf{/home/aholdengouveia/aholdengouveia.github.io/IntroLinux/labs/shellloop.html}}

The HTML version provides enhanced accessibility features including keyboard navigation, screen reader support, responsive design, dark mode support, and high contrast options.



\section*{Objectives:}
\begin{itemize}
    \item Learn how to do loops in bash scripting. Become comfortable and competent with the concept of count controlled and event controlled loops and how to implement them. Learn more about why testing and commenting is important.
\end{itemize}

\section*{Complete the following problems}

References, a video, a PowerPoint and some notes are available at my website
Link: \href{https://www.aholdengouveia.name/IntroLinux/shellloop.html}{\textbf{https://www.aholdengouveia.name/IntroLinux/shellloop.html}}
\subsection*{How to turn in Scripts}
    \begin{enumerate}
        \item Testing must be submitted along with your code. You should submit a copy of the code and a screenshot or the results of your test so that I can see a successful run of your code. Every script should be tested.  A script that doesn't run is worth 0 points.
        \item Comments should be done on EVERY line of code, (this is very bad practice in the real world, never do this for a job) so that I can see that you understand the code that you are submitting. Every line should have an explanation of what that line does.  Comments are worth 5 points per script.
        \item For submission you need to submit a copy of your code (actual text NO screenshots of code allowed) and your successful test run of your code for every script.
\end{enumerate}


\subsection*{Scripts}
    \begin{enumerate}
        \item Write a script that will echo a phrase for a given number of times. The user should generate the phrase and the number of times it is echoed. You may use read or positional parameters, but you need to include a comment about which one you picked and why.  Include a comment with your name.
        \item Write a script that will do a countdown. The user should be able to specify what number is started with, the countdown should go to 1.
    \end{enumerate}

It's better to start simple and work your way up.  I have a video on my Linux FAQ playlist on YouTube if you need help. Don't try and do the whole nested statements all at once, start with the basics, use my example if you want, and make small changes. 

\section*{Deliverables}
One file that includes all the requested information with your name, the date, and the lab title, make sure to follow the requested pattern for scripts, an example is included with the first scripting lab
\end{document}